% =====================================================================
\section{Компоненты для курса}
% =====================================================================
Список рассчитан на одно рабочее место. Детали из верхней части таблицы нужны почти в каждом уроке;
нижняя часть --- для отдельных опытов и заданий.

{\small
\begin{xltabular}{\textwidth}{@{}X l l@{}}
\toprule
Компонент & Кол-во & Уроки \\
\midrule
Плата ESP32-S3-DevKitC-1 (модуль N8R8 или N16R8) & 1 & все \\
USB-кабель \emph{с передачей данных} (Type-C или micro-USB по плате) & 1 & все \\
Макетная плата на 830 точек и набор проводов «папа--папа» & 1 & все \\
Светодиоды 5 мм: красный, жёлтый, зелёный & по 2 & 1--4, 7, 8 \\
Резисторы \SI{220}{\ohm} & 10 & 1--4 \\
Резисторы \SI{10}{k\ohm}, \SI{33}{k\ohm}, \SI{100}{k\ohm} & по 2 & 3, 5 \\
Конденсаторы \SI{100}{nF} & 5 & 5, 10 \\
Тактовые кнопки & 3 & 3, 7, 9, 40 \\
Потенциометр \SI{10}{k\ohm} & 1 & 5--7, 39 \\
OLED-дисплей SSD1306 128×64, интерфейс I\textsuperscript{2}C & 1 & 6, 7 \\
\midrule
Фоторезистор & 1 & 5 (задание) \\
RGB-светодиод с общим катодом & 1 & 4 (задание) \\
Пьезопищалка (пассивная) & 1 & 4 (задание), 24 \\
Сервопривод SG90 и отдельный источник \SI{5}{V} & 1 & 4 \\
Датчик BME280 или AHT20 (I\textsuperscript{2}C) & 1 & 6 (задание) \\
Модуль microSD с логикой 3,3 В и карта памяти & 1 & 10 \\
Вторая плата ESP32-S3 (или любая ESP32) & 1 & 2, 12 \\
Мультиметр & 1 & 5, 9, 14 \\
Датчик температуры DS18B20 (корпус или герметичный щуп) и резистор \SI{4.7}{k\ohm} & 1 & 15 \\
Датчик температуры и влажности DHT22 (AM2302), модуль с резистором & 1 & 15 \\
Ультразвуковой дальномер HC-SR04 и резисторы \SI{1}{k\ohm} и \SI{2}{k\ohm} для делителя & 1 & 15 \\
N-MOSFET логического уровня (IRLZ44N, AO3400 или аналог) & 2 & 16 \\
Диоды 1N4007 или 1N4148, резисторы \SI{100}{\ohm} и \SI{10}{k\ohm} & по 2 & 16 \\
Модуль реле на 1 канал (5 В, с оптопарой) & 1 & 16 \\
Драйвер мотора TB6612FNG или DRV8833 (модуль) & 1 & 16 \\
Коллекторный мотор постоянного тока \SIrange{3}{6}{V} (например, N20 или TT) & 1 & 16 \\
Лента, кольцо или модуль WS2812B на 8 пикселей, резистор \SI{330}{\ohm}, конденсатор \SI{1000}{\micro F} & 1 & 17 \\
Блок питания \SI{5}{V} на \SI{2}{A} и более (для ленты и мотора) & 1 & 16, 17 \\
Плата ESP32-S3-CAM с камерой OV2640, шлейф и microSD-карта & 1 & 20 \\
Светодиоды для светофора: красный, жёлтый, зелёный (из основного набора) & по 1 & 21 \\
Сдвиговый регистр 74HC595 и семисегментный индикатор с общим катодом & по 1 & 22 \\
Модуль энкодера KY-040 и матричная клавиатура 4×4 & по 1 & 23 \\
Шаговый двигатель 28BYJ-48 с платой драйвера ULN2003 & 1 & 25 \\
Модуль акселерометра и гироскопа MPU6050 (GY-521) & 1 & 26, 42 \\
Держатель на три батарейки AA или powerbank & 1 & 27 \\
Li-ion аккумулятор с платой защиты и модуль зарядки TP4056 (по решению педагога) & 1 & 27 \\
Цифровой микрофон INMP441 & 1 & 29 \\
Логический анализатор с входами 3,3 В & 1 на группу & 30, 31 \\
CAN-трансиверы с логикой 3,3 В и резисторы 120 Ом & 2 & 32 \\
RS-485-трансиверы с логикой 3,3 В & 2 & 32 \\
SPI-дисплей ST7789 240×240 с CS & 1 & 33 \\
Плата с RGB-экраном / USB Host и управляемым VBUS & по желанию & 33, 34 \\
Кнопка и внешний резистор 10 кОм & по 1 & 36 \\
\bottomrule
\end{xltabular}
}

\begin{tip}
Проверяйте, что USB-кабель передаёт данные: многие кабели от зарядных устройств только питают
плату, и тогда компьютер её «не видит».
\end{tip}
